\documentclass[11pt,a4paper]{article}

%=========================================================
% PACKAGES
%=========================================================
\usepackage[utf8]{inputenc}
\usepackage[T1]{fontenc}
\usepackage[french]{babel}
\usepackage[a4paper,margin=2.1cm]{geometry}
\usepackage{lmodern}
\usepackage{amsmath,amssymb,amsthm,mathtools}
\usepackage{fancyhdr}
\usepackage{lastpage}
\usepackage{enumitem}
\usepackage{array,tabularx}
\usepackage{tikz}
\usepackage[most]{tcolorbox}
\usepackage{hyperref}
\usepackage{xcolor}

%=========================================================
% MISE EN PAGE
%=========================================================
\setlength{\headheight}{14pt}
\pagestyle{fancy}
\fancyhf{}
\fancyhead[L]{Analyse CPGE}
\fancyhead[C]{Dérivation et fonctions usuelles}
\fancyhead[R]{Page \thepage/\pageref{LastPage}}
\fancyfoot[C]{Dérivabilité -- Rolle -- Accroissements finis -- Fonctions usuelles}

\hypersetup{
    colorlinks=true,
    linkcolor=blue!50!black,
    urlcolor=blue!50!black
}

\setlist[itemize]{label=\(\triangleright\)}
\setlist[enumerate]{label=\textbf{\arabic*.}}

%=========================================================
% COULEURS ET BOÎTES
%=========================================================
\definecolor{bleufonce}{RGB}{25,65,130}
\definecolor{bleuclair}{RGB}{235,243,255}
\definecolor{vertfonce}{RGB}{30,120,70}
\definecolor{vertclair}{RGB}{235,250,240}
\definecolor{rougefonce}{RGB}{170,40,40}
\definecolor{rougeclair}{RGB}{255,238,238}
\definecolor{orangefonce}{RGB}{190,110,20}
\definecolor{orangeclair}{RGB}{255,246,230}
\definecolor{violetfonce}{RGB}{95,55,145}
\definecolor{violetclair}{RGB}{246,241,255}

\tcbset{
    enhanced,
    sharp corners,
    boxrule=0.75pt,
    before skip=8pt,
    after skip=8pt
}

\newtcolorbox{definitionbox}[1][]{
    colback=bleuclair,
    colframe=bleufonce,
    title=\textbf{Définition #1}
}

\newtcolorbox{theorembox}[1][]{
    colback=vertclair,
    colframe=vertfonce,
    title=\textbf{Théorème / Propriété #1}
}

\newtcolorbox{methodbox}[1][]{
    colback=orangeclair,
    colframe=orangefonce,
    title=\textbf{Méthode #1}
}

\newtcolorbox{retenirbox}[1][]{
    colback=violetclair,
    colframe=violetfonce,
    title=\textbf{À retenir #1}
}

\newtcolorbox{erreurbox}[1][]{
    colback=rougeclair,
    colframe=rougefonce,
    title=\textbf{Erreur classique #1}
}

\newtcolorbox{exercicebox}[1][]{
    colback=white,
    colframe=black!60,
    title=\textbf{Exercice #1}
}

\newtcolorbox{correctionbox}[1][]{
    colback=gray!5,
    colframe=black!65,
    title=\textbf{Correction #1}
}

%=========================================================
% COMMANDES
%=========================================================
\newcommand{\N}{\mathbb N}
\newcommand{\Z}{\mathbb Z}
\newcommand{\Q}{\mathbb Q}
\newcommand{\R}{\mathbb R}
\newcommand{\C}{\mathbb C}
\newcommand{\abs}[1]{\left|#1\right|}
\newcommand{\ssi}{\Longleftrightarrow}
\newcommand{\implique}{\Longrightarrow}
\newcommand{\e}{\mathrm e}
\newcommand{\dd}{\mathrm d}
\newcommand{\Id}{\operatorname{Id}}
\newcommand{\arccot}{\operatorname{arccot}}
\newcommand{\ch}{\operatorname{ch}}
\newcommand{\sh}{\operatorname{sh}}
\renewcommand{\th}{\operatorname{th}}

\newtheorem{prop}{Proposition}[section]
\newtheorem{lemme}[prop]{Lemme}
\newtheorem{corollaire}[prop]{Corollaire}

%=========================================================
% DOCUMENT
%=========================================================
\begin{document}

\begin{center}
    {\Huge \textbf{Dérivation, théorèmes fondamentaux et fonctions usuelles}}\\[2mm]
    {\Large Chapitre 4 -- Analyse réelle en classe préparatoire scientifique}\\[2mm]
    {\large Cours complet, progressif et rigoureux}\\[2mm]
    \rule{0.88\linewidth}{0.5pt}
\end{center}

\vspace{0.4cm}

\tableofcontents

\newpage

%=========================================================
\section*{Introduction générale}
\addcontentsline{toc}{section}{Introduction générale}

La dérivation est l'un des outils les plus puissants de l'analyse réelle.  
Elle permet d'étudier localement une fonction, de mesurer sa variation instantanée, de déterminer ses extrema, de prouver des inégalités, d'étudier les variations et de construire rigoureusement les fonctions usuelles.

Dans ce chapitre, on part de la définition locale de la dérivée en un point, puis on établit les règles de calcul.  
On démontre ensuite les grands théorèmes fondamentaux : le théorème de Rolle, le théorème des accroissements finis et l'inégalité des accroissements finis.  
Enfin, on étudie les fonctions usuelles : logarithme, exponentielle, puissances, fonctions trigonométriques, fonctions trigonométriques réciproques et fonctions hyperboliques.

\begin{retenirbox}
La dérivée est d'abord une notion locale :
\[
f'(a)=\lim_{x\to a}\frac{f(x)-f(a)}{x-a}.
\]
Elle devient ensuite un outil global grâce aux théorèmes de Rolle et des accroissements finis.
\end{retenirbox}

%=========================================================
\section{Dérivabilité en un point}

\subsection{Taux d'accroissement}

\begin{definitionbox}[Taux d'accroissement]
Soit \(f:I\to\R\), où \(I\) est un intervalle de \(\R\).  
Soient \(a,x\in I\), avec \(x\neq a\).  
Le taux d'accroissement de \(f\) entre \(a\) et \(x\) est :
\[
\frac{f(x)-f(a)}{x-a}.
\]
\end{definitionbox}

Ce quotient mesure la pente de la droite sécante au graphe de \(f\) passant par les points :
\[
(a,f(a))
\quad\text{et}\quad
(x,f(x)).
\]

\begin{retenirbox}
Le taux d'accroissement est une pente moyenne.  
La dérivée est la limite de cette pente moyenne lorsque \(x\) tend vers \(a\).
\end{retenirbox}

\subsection{Définition de la dérivabilité}

\begin{definitionbox}[Dérivabilité en un point]
Soit \(f:I\to\R\), où \(I\) est un intervalle de \(\R\).  
Soit \(a\in I\).  
On dit que \(f\) est dérivable en \(a\) si la limite :
\[
\lim_{x\to a,\ x\neq a}\frac{f(x)-f(a)}{x-a}
\]
existe et est finie.

Cette limite est appelée nombre dérivé de \(f\) en \(a\), et est notée :
\[
f'(a).
\]
\end{definitionbox}

On peut aussi écrire, en posant \(h=x-a\) :
\[
f'(a)=\lim_{h\to0}\frac{f(a+h)-f(a)}{h}.
\]

\begin{retenirbox}
Deux écritures équivalentes :
\[
f'(a)=\lim_{x\to a}\frac{f(x)-f(a)}{x-a},
\]
et :
\[
f'(a)=\lim_{h\to0}\frac{f(a+h)-f(a)}{h}.
\]
\end{retenirbox}

\subsection{Dérivabilité à gauche et à droite}

\begin{definitionbox}[Dérivée à droite et à gauche]
Soit \(f:I\to\R\) et \(a\in I\).  
Si la limite :
\[
\lim_{h\to0^+}\frac{f(a+h)-f(a)}{h}
\]
existe, elle est appelée dérivée à droite de \(f\) en \(a\), notée :
\[
f'_d(a).
\]

Si la limite :
\[
\lim_{h\to0^-}\frac{f(a+h)-f(a)}{h}
\]
existe, elle est appelée dérivée à gauche de \(f\) en \(a\), notée :
\[
f'_g(a).
\]
\end{definitionbox}

\begin{theorembox}
La fonction \(f\) est dérivable en \(a\) si et seulement si \(f\) est dérivable à gauche et à droite en \(a\), et si :
\[
f'_g(a)=f'_d(a).
\]
Dans ce cas :
\[
f'(a)=f'_g(a)=f'_d(a).
\]
\end{theorembox}

\begin{exercicebox}[Dérivée à gauche et à droite]
Étudier la dérivabilité en \(0\) de :
\[
f(x)=|x|.
\]
\end{exercicebox}

\begin{correctionbox}
On calcule le taux d'accroissement en \(0\) :
\[
\frac{f(h)-f(0)}{h}=\frac{|h|}{h}.
\]
Si \(h>0\), alors :
\[
\frac{|h|}{h}=1.
\]
Donc :
\[
f'_d(0)=1.
\]
Si \(h<0\), alors :
\[
\frac{|h|}{h}=-1.
\]
Donc :
\[
f'_g(0)=-1.
\]
Les deux dérivées latérales sont différentes.  
Ainsi, \(f\) n'est pas dérivable en \(0\).
\end{correctionbox}

%=========================================================
\section{Dérivée comme approximation locale}

\subsection{Approximation affine}

\begin{theorembox}[Dérivabilité et développement limité d'ordre 1]
Soit \(f:I\to\R\) et \(a\in I\).  
La fonction \(f\) est dérivable en \(a\) si et seulement s'il existe un réel \(L\) et une fonction \(\varepsilon\) définie au voisinage de \(0\), avec :
\[
\varepsilon(h)\to0
\quad\text{quand}\quad h\to0,
\]
tels que :
\[
f(a+h)=f(a)+Lh+h\varepsilon(h).
\]
Dans ce cas :
\[
L=f'(a).
\]
\end{theorembox}

\begin{proof}
Supposons que \(f\) soit dérivable en \(a\).  
Pour \(h\neq0\), on écrit :
\[
\frac{f(a+h)-f(a)}{h}=f'(a)+\varepsilon(h),
\]
où :
\[
\varepsilon(h)=\frac{f(a+h)-f(a)}{h}-f'(a).
\]
Par définition de la dérivée :
\[
\varepsilon(h)\to0.
\]
Donc :
\[
f(a+h)=f(a)+hf'(a)+h\varepsilon(h).
\]

Réciproquement, supposons :
\[
f(a+h)=f(a)+Lh+h\varepsilon(h),
\]
avec :
\[
\varepsilon(h)\to0.
\]
Alors, pour \(h\neq0\),
\[
\frac{f(a+h)-f(a)}{h}=L+\varepsilon(h).
\]
En faisant tendre \(h\) vers \(0\), on obtient :
\[
\frac{f(a+h)-f(a)}{h}\to L.
\]
Donc \(f\) est dérivable en \(a\), et :
\[
f'(a)=L.
\]
\end{proof}

\begin{retenirbox}
La dérivabilité en \(a\) signifie que localement :
\[
f(a+h)=f(a)+hf'(a)+o(h).
\]
La fonction est donc bien approchée, au premier ordre, par l'application affine :
\[
x\mapsto f(a)+f'(a)(x-a).
\]
\end{retenirbox}

\subsection{Tangente au graphe}

Si \(f\) est dérivable en \(a\), la tangente au graphe de \(f\) au point d'abscisse \(a\) a pour équation :
\[
y=f(a)+f'(a)(x-a).
\]

\begin{methodbox}[Trouver une équation de tangente]
Pour déterminer la tangente à la courbe de \(f\) au point d'abscisse \(a\) :
\begin{enumerate}
    \item calculer \(f(a)\) ;
    \item calculer \(f'(a)\) ;
    \item écrire :
    \[
    y=f(a)+f'(a)(x-a).
    \]
\end{enumerate}
\end{methodbox}

\begin{exercicebox}[Tangente]
Déterminer la tangente à la courbe de :
\[
f(x)=x^2
\]
au point d'abscisse \(a=1\).
\end{exercicebox}

\begin{correctionbox}
On a :
\[
f(1)=1.
\]
De plus :
\[
f'(x)=2x,
\]
donc :
\[
f'(1)=2.
\]
L'équation de la tangente est :
\[
y=1+2(x-1),
\]
c'est-à-dire :
\[
y=2x-1.
\]
\end{correctionbox}

%=========================================================
\section{Dérivabilité et continuité}

\begin{theorembox}
Si \(f\) est dérivable en \(a\), alors \(f\) est continue en \(a\).
\end{theorembox}

\begin{proof}
Si \(f\) est dérivable en \(a\), alors :
\[
f(a+h)=f(a)+hf'(a)+h\varepsilon(h),
\]
avec :
\[
\varepsilon(h)\to0.
\]
Donc :
\[
f(a+h)-f(a)=hf'(a)+h\varepsilon(h).
\]
Le membre de droite tend vers \(0\) lorsque \(h\to0\).  
Ainsi :
\[
f(a+h)\to f(a).
\]
Donc \(f\) est continue en \(a\).
\end{proof}

\begin{erreurbox}
La réciproque est fausse.  
Une fonction peut être continue sans être dérivable.  
Exemple :
\[
x\mapsto |x|
\]
est continue en \(0\), mais n'est pas dérivable en \(0\).
\end{erreurbox}

%=========================================================
\section{Fonction dérivée et classe \(\mathcal C^1\)}

\begin{definitionbox}[Dérivabilité sur un intervalle]
Soit \(f:I\to\R\).  
On dit que \(f\) est dérivable sur \(I\) si elle est dérivable en tout point de \(I\), en tenant compte des dérivées latérales aux bornes éventuelles de \(I\).
\end{definitionbox}

\begin{definitionbox}[Fonction dérivée]
Si \(f\) est dérivable sur \(I\), on définit la fonction dérivée :
\[
f':I\to\R,
\qquad
x\mapsto f'(x).
\]
\end{definitionbox}

\begin{definitionbox}[Fonction de classe \(\mathcal C^1\)]
On dit que \(f\) est de classe \(\mathcal C^1\) sur \(I\) si \(f\) est dérivable sur \(I\) et si sa dérivée \(f'\) est continue sur \(I\).
\end{definitionbox}

\begin{retenirbox}
Dérivable ne signifie pas nécessairement de classe \(\mathcal C^1\).  
La classe \(\mathcal C^1\) demande en plus la continuité de la dérivée.
\end{retenirbox}

%=========================================================
\section{Opérations sur les dérivées}

\subsection{Linéarité}

\begin{theorembox}[Linéarité de la dérivation]
Soient \(f,g:I\to\R\) dérivables en \(a\), et soient \(\lambda,\mu\in\R\).  
Alors :
\[
\lambda f+\mu g
\]
est dérivable en \(a\), et :
\[
(\lambda f+\mu g)'(a)=\lambda f'(a)+\mu g'(a).
\]
\end{theorembox}

\begin{proof}
On écrit :
\[
\frac{(\lambda f+\mu g)(a+h)-(\lambda f+\mu g)(a)}{h}
=
\lambda\frac{f(a+h)-f(a)}{h}
+
\mu\frac{g(a+h)-g(a)}{h}.
\]
En faisant tendre \(h\) vers \(0\), on obtient :
\[
(\lambda f+\mu g)'(a)=\lambda f'(a)+\mu g'(a).
\]
\end{proof}

\subsection{Produit}

\begin{theorembox}[Dérivée d'un produit]
Si \(f\) et \(g\) sont dérivables en \(a\), alors \(fg\) est dérivable en \(a\), et :
\[
(fg)'(a)=f'(a)g(a)+f(a)g'(a).
\]
\end{theorembox}

\begin{proof}
On écrit :
\[
f(a+h)g(a+h)-f(a)g(a)
\]
sous la forme :
\[
\bigl(f(a+h)-f(a)\bigr)g(a+h)
+
f(a)\bigl(g(a+h)-g(a)\bigr).
\]
En divisant par \(h\), on obtient :
\[
\frac{f(a+h)g(a+h)-f(a)g(a)}{h}
=
\frac{f(a+h)-f(a)}{h}g(a+h)
+
f(a)\frac{g(a+h)-g(a)}{h}.
\]
Comme \(g\) est dérivable en \(a\), elle est continue en \(a\), donc :
\[
g(a+h)\to g(a).
\]
En faisant tendre \(h\) vers \(0\), on obtient :
\[
(fg)'(a)=f'(a)g(a)+f(a)g'(a).
\]
\end{proof}

\subsection{Inverse et quotient}

\begin{theorembox}[Dérivée de l'inverse]
Si \(g\) est dérivable en \(a\) et si :
\[
g(a)\neq0,
\]
alors \(\frac1g\) est dérivable en \(a\), et :
\[
\left(\frac1g\right)'(a)=-\frac{g'(a)}{g(a)^2}.
\]
\end{theorembox}

\begin{proof}
On calcule :
\[
\frac{\frac1{g(a+h)}-\frac1{g(a)}}{h}
=
\frac{g(a)-g(a+h)}{h\,g(a+h)g(a)}
=
-\frac{g(a+h)-g(a)}{h}\cdot\frac1{g(a+h)g(a)}.
\]
Comme \(g(a+h)\to g(a)\neq0\), on obtient :
\[
\left(\frac1g\right)'(a)=-\frac{g'(a)}{g(a)^2}.
\]
\end{proof}

\begin{theorembox}[Dérivée d'un quotient]
Si \(f\) et \(g\) sont dérivables en \(a\), et si \(g(a)\neq0\), alors :
\[
\frac fg
\]
est dérivable en \(a\), et :
\[
\left(\frac fg\right)'(a)
=
\frac{f'(a)g(a)-f(a)g'(a)}{g(a)^2}.
\]
\end{theorembox}

\begin{proof}
On écrit :
\[
\frac fg=f\cdot\frac1g.
\]
On applique la formule du produit et celle de l'inverse.
\end{proof}

\begin{methodbox}[Choisir la bonne règle de dérivation]
Avant de dériver, identifier la structure :
\begin{itemize}
    \item somme : utiliser la linéarité ;
    \item produit : utiliser \((fg)'=f'g+fg'\) ;
    \item quotient : utiliser \(\left(\frac fg\right)'=\frac{f'g-fg'}{g^2}\) ;
    \item composée : utiliser \((f\circ g)'=(f'\circ g)g'\).
\end{itemize}
\end{methodbox}

\begin{exercicebox}[Calcul de dérivée]
Dériver :
\[
f(x)=\frac{x^2+1}{x-1}.
\]
\end{exercicebox}

\begin{correctionbox}
On pose :
\[
u(x)=x^2+1,
\qquad
v(x)=x-1.
\]
Alors :
\[
u'(x)=2x,
\qquad
v'(x)=1.
\]
Sur \(\R\setminus\{1\}\),
\[
f'(x)=\frac{u'(x)v(x)-u(x)v'(x)}{v(x)^2}.
\]
Donc :
\[
f'(x)=\frac{2x(x-1)-(x^2+1)}{(x-1)^2}.
\]
On simplifie :
\[
f'(x)=\frac{2x^2-2x-x^2-1}{(x-1)^2}
=
\frac{x^2-2x-1}{(x-1)^2}.
\]
\end{correctionbox}

%=========================================================
\section{Dérivée d'une composée}

\begin{theorembox}[Dérivée d'une composée]
Soient \(g:I\to J\) et \(f:J\to\R\).  
Si \(g\) est dérivable en \(a\), et si \(f\) est dérivable en \(g(a)\), alors :
\[
f\circ g
\]
est dérivable en \(a\), et :
\[
(f\circ g)'(a)=f'(g(a))g'(a).
\]
\end{theorembox}

\begin{proof}
Comme \(f\) est dérivable en \(g(a)\), on peut écrire :
\[
f(g(a)+k)=f(g(a))+kf'(g(a))+k\varepsilon(k),
\]
avec :
\[
\varepsilon(k)\to0
\quad\text{quand}\quad k\to0.
\]
On pose :
\[
k=g(a+h)-g(a).
\]
Comme \(g\) est dérivable en \(a\), elle est continue en \(a\), donc :
\[
k\to0.
\]
Ainsi :
\[
f(g(a+h))-f(g(a))
=
\bigl(g(a+h)-g(a)\bigr)f'(g(a))
+
\bigl(g(a+h)-g(a)\bigr)\varepsilon(g(a+h)-g(a)).
\]
On divise par \(h\) :
\[
\frac{f(g(a+h))-f(g(a))}{h}
=
\frac{g(a+h)-g(a)}{h}f'(g(a))
+
\frac{g(a+h)-g(a)}{h}\varepsilon(g(a+h)-g(a)).
\]
En faisant tendre \(h\) vers \(0\), on obtient :
\[
(f\circ g)'(a)=g'(a)f'(g(a)).
\]
\end{proof}

\begin{retenirbox}
Formule fondamentale :
\[
(f\circ g)'=(f'\circ g)\,g'.
\]
Exemple :
\[
(\sin(x^2))'=2x\cos(x^2).
\]
\end{retenirbox}

\begin{exercicebox}[Composition]
Dériver :
\[
f(x)=\sqrt{1+x^2}.
\]
\end{exercicebox}

\begin{correctionbox}
On écrit :
\[
f(x)=(1+x^2)^{1/2}.
\]
C'est une composée de :
\[
u(x)=1+x^2
\]
et :
\[
\varphi(t)=\sqrt t.
\]
On a :
\[
u'(x)=2x,
\]
et :
\[
\varphi'(t)=\frac1{2\sqrt t}.
\]
Donc :
\[
f'(x)=\frac1{2\sqrt{1+x^2}}\cdot 2x
=
\frac{x}{\sqrt{1+x^2}}.
\]
\end{correctionbox}

%=========================================================
\section{Dérivée d'une fonction réciproque}

\subsection{Théorème général}

\begin{theorembox}[Dérivée d'une réciproque]
Soit \(f:I\to J\) une bijection continue strictement monotone.  
On suppose que \(f\) est dérivable en \(a\in I\), et que :
\[
f'(a)\neq0.
\]
Alors \(f^{-1}\) est dérivable en :
\[
b=f(a),
\]
et :
\[
(f^{-1})'(b)=\frac1{f'(a)}.
\]
Autrement dit :
\[
(f^{-1})'(y)=\frac1{f'(f^{-1}(y))}
\]
là où cette expression a un sens.
\end{theorembox}

\begin{proof}
Posons :
\[
b=f(a).
\]
Pour \(y\neq b\), on écrit :
\[
\frac{f^{-1}(y)-f^{-1}(b)}{y-b}.
\]
Posons :
\[
x=f^{-1}(y).
\]
Alors :
\[
y=f(x),
\]
et :
\[
f^{-1}(b)=a.
\]
Le quotient devient :
\[
\frac{x-a}{f(x)-f(a)}
=
\frac1{\frac{f(x)-f(a)}{x-a}}.
\]
Lorsque \(y\to b\), comme \(f^{-1}\) est continue, on a :
\[
x=f^{-1}(y)\to f^{-1}(b)=a.
\]
Donc :
\[
\frac{f(x)-f(a)}{x-a}\to f'(a).
\]
Comme \(f'(a)\neq0\), on obtient :
\[
\frac{x-a}{f(x)-f(a)}\to\frac1{f'(a)}.
\]
Ainsi :
\[
(f^{-1})'(b)=\frac1{f'(a)}.
\]
\end{proof}

\begin{retenirbox}
La dérivée de la réciproque est l'inverse de la dérivée, mais au point correspondant :
\[
(f^{-1})'(y)=\frac1{f'(f^{-1}(y))}.
\]
\end{retenirbox}

\begin{erreurbox}
Ne pas écrire :
\[
(f^{-1})'(x)=\frac1{f'(x)}.
\]
La bonne formule est :
\[
(f^{-1})'(x)=\frac1{f'(f^{-1}(x))}.
\]
\end{erreurbox}

%=========================================================
\section{Extrema locaux et points critiques}

\subsection{Définitions}

\begin{definitionbox}[Maximum local]
Soit \(f:I\to\R\) et \(a\in I\).  
On dit que \(f\) admet un maximum local en \(a\) s'il existe \(\eta>0\) tel que :
\[
\forall x\in I,\quad |x-a|\leq\eta\implique f(x)\leq f(a).
\]
\end{definitionbox}

\begin{definitionbox}[Minimum local]
On dit que \(f\) admet un minimum local en \(a\) s'il existe \(\eta>0\) tel que :
\[
\forall x\in I,\quad |x-a|\leq\eta\implique f(a)\leq f(x).
\]
\end{definitionbox}

\begin{definitionbox}[Point critique]
Si \(f\) est dérivable en \(a\), on dit que \(a\) est un point critique de \(f\) si :
\[
f'(a)=0.
\]
\end{definitionbox}

\subsection{Condition nécessaire d'extremum local}

\begin{theorembox}[Condition nécessaire d'extremum local]
Soit \(f:I\to\R\) dérivable en un point \(a\) intérieur à \(I\).  
Si \(f\) admet un extremum local en \(a\), alors :
\[
f'(a)=0.
\]
\end{theorembox}

\begin{proof}
Supposons que \(f\) admette un maximum local en \(a\).  
Il existe \(\eta>0\) tel que, pour \(x\) proche de \(a\),
\[
f(x)\leq f(a).
\]

Pour \(h>0\) suffisamment petit :
\[
\frac{f(a+h)-f(a)}{h}\leq0.
\]
En faisant tendre \(h\to0^+\), on obtient :
\[
f'_d(a)\leq0.
\]

Pour \(h<0\) suffisamment proche de \(0\), on a encore :
\[
f(a+h)-f(a)\leq0.
\]
Mais comme \(h<0\), en divisant par \(h\), le sens de l'inégalité change :
\[
\frac{f(a+h)-f(a)}{h}\geq0.
\]
En faisant tendre \(h\to0^-\), on obtient :
\[
f'_g(a)\geq0.
\]

Comme \(f\) est dérivable en \(a\), les dérivées latérales sont égales :
\[
f'_g(a)=f'_d(a)=f'(a).
\]
Donc :
\[
f'(a)\geq0
\quad\text{et}\quad
f'(a)\leq0.
\]
Ainsi :
\[
f'(a)=0.
\]

Le cas d'un minimum local est analogue.
\end{proof}

\begin{erreurbox}
Un point critique n'est pas forcément un extremum local.  
Exemple :
\[
f(x)=x^3.
\]
On a :
\[
f'(0)=0,
\]
mais \(0\) n'est ni un maximum local, ni un minimum local.
\end{erreurbox}

%=========================================================
\section{Théorème de Rolle}

\begin{theorembox}[Théorème de Rolle]
Soit \(f:[a,b]\to\R\) une fonction continue sur \([a,b]\), dérivable sur \(]a,b[\), avec :
\[
f(a)=f(b).
\]
Alors il existe \(c\in]a,b[\) tel que :
\[
f'(c)=0.
\]
\end{theorembox}

\begin{proof}
Comme \(f\) est continue sur le segment \([a,b]\), elle atteint son maximum et son minimum.  
Il existe donc \(\alpha,\beta\in[a,b]\) tels que :
\[
f(\alpha)=\min_{[a,b]}f,
\qquad
f(\beta)=\max_{[a,b]}f.
\]

Si \(f\) est constante sur \([a,b]\), alors :
\[
f'(x)=0
\]
pour tout \(x\in]a,b[\), et le résultat est immédiat.

Sinon, le minimum et le maximum ne sont pas égaux.  
Comme :
\[
f(a)=f(b),
\]
au moins l'un des deux extrema est atteint en un point intérieur à \(]a,b[\).  
Appelons ce point \(c\).  
La fonction \(f\) admet donc un extremum local en \(c\), et \(c\) est intérieur à l'intervalle.  
Par la condition nécessaire d'extremum local :
\[
f'(c)=0.
\]
\end{proof}

\begin{methodbox}[Utiliser Rolle]
Pour appliquer le théorème de Rolle, vérifier :
\begin{enumerate}
    \item continuité sur le segment \([a,b]\) ;
    \item dérivabilité sur l'intervalle ouvert \(]a,b[\) ;
    \item égalité des valeurs aux extrémités :
    \[
    f(a)=f(b).
    \]
\end{enumerate}
\end{methodbox}

\begin{exercicebox}[Application de Rolle]
Montrer que si un polynôme réel \(P\) possède deux racines réelles distinctes, alors \(P'\) possède au moins une racine réelle entre elles.
\end{exercicebox}

\begin{correctionbox}
Soient \(a<b\) deux racines de \(P\).  
Alors :
\[
P(a)=P(b)=0.
\]
Un polynôme est continu sur \([a,b]\) et dérivable sur \(]a,b[\).  
Par le théorème de Rolle, il existe \(c\in]a,b[\) tel que :
\[
P'(c)=0.
\]
Donc \(P'\) possède au moins une racine réelle dans \(]a,b[\).
\end{correctionbox}

%=========================================================
\section{Théorème des accroissements finis}

\subsection{Énoncé}

\begin{theorembox}[Théorème des accroissements finis]
Soit \(f:[a,b]\to\R\) une fonction continue sur \([a,b]\) et dérivable sur \(]a,b[\).  
Alors il existe \(c\in]a,b[\) tel que :
\[
f(b)-f(a)=f'(c)(b-a).
\]
Autrement dit :
\[
\frac{f(b)-f(a)}{b-a}=f'(c).
\]
\end{theorembox}

\begin{proof}
On considère la fonction auxiliaire :
\[
g(x)=f(x)-\frac{f(b)-f(a)}{b-a}(x-a).
\]
La fonction \(g\) est continue sur \([a,b]\) et dérivable sur \(]a,b[\).  
On calcule :
\[
g(a)=f(a),
\]
et :
\[
g(b)=f(b)-\frac{f(b)-f(a)}{b-a}(b-a)=f(a).
\]
Ainsi :
\[
g(a)=g(b).
\]
Par le théorème de Rolle, il existe \(c\in]a,b[\) tel que :
\[
g'(c)=0.
\]
Or :
\[
g'(x)=f'(x)-\frac{f(b)-f(a)}{b-a}.
\]
Donc :
\[
f'(c)=\frac{f(b)-f(a)}{b-a}.
\]
\end{proof}

\begin{retenirbox}
Le théorème des accroissements finis affirme que, pour une fonction dérivable, une pente moyenne est égale à une pente instantanée.
\[
\frac{f(b)-f(a)}{b-a}=f'(c).
\]
\end{retenirbox}

\subsection{Applications aux variations}

\begin{theorembox}[Lien entre signe de la dérivée et variations]
Soit \(f:I\to\R\) dérivable sur un intervalle \(I\).
\begin{enumerate}
    \item Si :
    \[
    \forall x\in I,\quad f'(x)\geq0,
    \]
    alors \(f\) est croissante sur \(I\).
    \item Si :
    \[
    \forall x\in I,\quad f'(x)\leq0,
    \]
    alors \(f\) est décroissante sur \(I\).
    \item Si :
    \[
    \forall x\in I,\quad f'(x)>0,
    \]
    alors \(f\) est strictement croissante sur \(I\).
    \item Si :
    \[
    \forall x\in I,\quad f'(x)<0,
    \]
    alors \(f\) est strictement décroissante sur \(I\).
\end{enumerate}
\end{theorembox}

\begin{proof}
Soient \(x<y\) dans \(I\).  
Par le théorème des accroissements finis appliqué sur \([x,y]\), il existe \(c\in]x,y[\) tel que :
\[
f(y)-f(x)=f'(c)(y-x).
\]
Comme :
\[
y-x>0,
\]
le signe de \(f(y)-f(x)\) est celui de \(f'(c)\).  
On obtient donc les conclusions annoncées.
\end{proof}

\begin{erreurbox}
La condition \(f'(x)>0\) pour tout \(x\) implique que \(f\) est strictement croissante.  
Mais une fonction peut être strictement croissante même si sa dérivée s'annule en certains points.

Exemple :
\[
f(x)=x^3
\]
est strictement croissante sur \(\R\), mais :
\[
f'(0)=0.
\]
\end{erreurbox}

%=========================================================
\section{Inégalité des accroissements finis}

\begin{theorembox}[Inégalité des accroissements finis]
Soit \(f:I\to\R\) dérivable sur un intervalle \(I\).  
On suppose qu'il existe \(K\geq0\) tel que :
\[
\forall x\in I,\quad |f'(x)|\leq K.
\]
Alors :
\[
\forall x,y\in I,\quad |f(x)-f(y)|\leq K|x-y|.
\]
\end{theorembox}

\begin{proof}
Soient \(x,y\in I\), avec \(x<y\).  
Par le théorème des accroissements finis appliqué sur \([x,y]\), il existe \(c\in]x,y[\) tel que :
\[
f(y)-f(x)=f'(c)(y-x).
\]
Donc :
\[
|f(y)-f(x)|=|f'(c)||y-x|.
\]
Comme :
\[
|f'(c)|\leq K,
\]
on obtient :
\[
|f(y)-f(x)|\leq K|y-x|.
\]
Le cas \(x=y\) est évident.
\end{proof}

\begin{definitionbox}[Fonction lipschitzienne]
Une fonction \(f:I\to\R\) est dite \(K\)-lipschitzienne si :
\[
\forall x,y\in I,\quad |f(x)-f(y)|\leq K|x-y|.
\]
\end{definitionbox}

\begin{retenirbox}
Si une fonction dérivable a une dérivée bornée par \(K\), alors elle est \(K\)-lipschitzienne.
\end{retenirbox}

\begin{exercicebox}[Inégalité des accroissements finis]
Montrer que :
\[
\forall x,y\in\R,\quad |\sin x-\sin y|\leq |x-y|.
\]
\end{exercicebox}

\begin{correctionbox}
La fonction sinus est dérivable sur \(\R\), et :
\[
(\sin)'=\cos.
\]
Or :
\[
|\cos t|\leq1
\]
pour tout \(t\in\R\).  
Par l'inégalité des accroissements finis avec \(K=1\), on obtient :
\[
|\sin x-\sin y|\leq |x-y|.
\]
\end{correctionbox}

%=========================================================
\section{Fonction logarithme népérien}

\subsection{Définition et propriétés}

\begin{definitionbox}[Logarithme népérien]
La fonction logarithme népérien, notée \(\ln\), est définie sur \(]0,+\infty[\).  
Elle vérifie :
\[
\ln(1)=0,
\]
\[
\forall x,y>0,\quad \ln(xy)=\ln x+\ln y,
\]
et :
\[
\ln'(x)=\frac1x.
\]
\end{definitionbox}

\begin{theorembox}[Propriétés algébriques]
Pour tous \(x,y>0\) et tout \(r\in\Q\) lorsque l'expression a un sens :
\[
\ln(xy)=\ln x+\ln y,
\]
\[
\ln\left(\frac{x}{y}\right)=\ln x-\ln y,
\]
\[
\ln(x^r)=r\ln x.
\]
\end{theorembox}

\begin{theorembox}[Variations et limites]
La fonction \(\ln\) est strictement croissante sur \(]0,+\infty[\).  
De plus :
\[
\lim_{x\to0^+}\ln x=-\infty,
\]
et :
\[
\lim_{x\to+\infty}\ln x=+\infty.
\]
\end{theorembox}

\begin{proof}
Comme :
\[
\ln'(x)=\frac1x>0
\]
sur \(]0,+\infty[\), la fonction \(\ln\) est strictement croissante.

Les limites aux bornes sont des résultats fondamentaux de la fonction logarithme.  
Elles sont cohérentes avec la croissance stricte et la propriété :
\[
\ln(x^n)=n\ln x.
\]
\end{proof}

\begin{methodbox}[Dériver une expression avec logarithme]
Pour dériver :
\[
\ln(u(x)),
\]
on utilise la formule de composition :
\[
(\ln u)'=\frac{u'}{u},
\]
sur les intervalles où :
\[
u(x)>0.
\]
\end{methodbox}

\begin{exercicebox}[Logarithme]
Dériver :
\[
f(x)=\ln(1+x^2).
\]
\end{exercicebox}

\begin{correctionbox}
On pose :
\[
u(x)=1+x^2.
\]
Alors :
\[
u'(x)=2x,
\]
et :
\[
u(x)>0
\]
pour tout \(x\in\R\).  
Donc :
\[
f'(x)=\frac{u'(x)}{u(x)}
=
\frac{2x}{1+x^2}.
\]
\end{correctionbox}

%=========================================================
\section{Fonction exponentielle}

\subsection{Définition}

\begin{definitionbox}[Exponentielle]
La fonction exponentielle est la fonction réciproque du logarithme népérien :
\[
\exp:\R\to]0,+\infty[.
\]
On note :
\[
\exp(x)=\e^x.
\]
\end{definitionbox}

On a donc :
\[
\ln(\e^x)=x
\quad\text{pour tout }x\in\R,
\]
et :
\[
\e^{\ln x}=x
\quad\text{pour tout }x>0.
\]

\begin{theorembox}[Propriétés de l'exponentielle]
Pour tous \(x,y\in\R\) :
\[
\e^{x+y}=\e^x\e^y,
\]
\[
\e^{-x}=\frac1{\e^x},
\]
\[
\e^0=1.
\]
\end{theorembox}

\begin{theorembox}[Dérivée de l'exponentielle]
La fonction exponentielle est dérivable sur \(\R\), et :
\[
(\e^x)'=\e^x.
\]
Plus généralement :
\[
(\e^{u(x)})'=u'(x)\e^{u(x)}.
\]
\end{theorembox}

\begin{proof}
Comme \(\exp\) est la réciproque de \(\ln\), on utilise la formule de dérivation d'une réciproque.  
Soit :
\[
y=\e^x.
\]
Alors :
\[
x=\ln y.
\]
Comme :
\[
(\ln)'(y)=\frac1y,
\]
on obtient :
\[
(\exp)'(x)=\frac1{(\ln)'(\exp x)}
=
\frac1{1/\e^x}
=
\e^x.
\]
\end{proof}

\begin{theorembox}[Variations et limites]
La fonction exponentielle est strictement croissante sur \(\R\).  
De plus :
\[
\lim_{x\to-\infty}\e^x=0,
\]
et :
\[
\lim_{x\to+\infty}\e^x=+\infty.
\]
\end{theorembox}

\begin{exercicebox}[Exponentielle]
Dériver :
\[
f(x)=\e^{x^2-3x}.
\]
\end{exercicebox}

\begin{correctionbox}
On pose :
\[
u(x)=x^2-3x.
\]
Alors :
\[
u'(x)=2x-3.
\]
Donc :
\[
f'(x)=(2x-3)\e^{x^2-3x}.
\]
\end{correctionbox}

%=========================================================
\section{Fonctions puissances}

\subsection{Puissances entières}

Pour \(n\in\N\), la fonction :
\[
x\mapsto x^n
\]
est définie sur \(\R\) et dérivable sur \(\R\), avec :
\[
(x^n)'=nx^{n-1}
\]
pour \(n\geq1\).

\subsection{Puissances réelles}

\begin{definitionbox}[Puissance réelle]
Pour \(a\in\R\) et \(x>0\), on définit :
\[
x^a=\e^{a\ln x}.
\]
\end{definitionbox}

\begin{theorembox}[Dérivée d'une puissance réelle]
Pour \(a\in\R\), la fonction :
\[
x\mapsto x^a
\]
est dérivable sur \(]0,+\infty[\), et :
\[
(x^a)'=a x^{a-1}.
\]
\end{theorembox}

\begin{proof}
On écrit :
\[
x^a=\e^{a\ln x}.
\]
Par dérivation d'une composée :
\[
(x^a)'=\frac{a}{x}\e^{a\ln x}.
\]
Donc :
\[
(x^a)'=\frac{a}{x}x^a=ax^{a-1}.
\]
\end{proof}

\begin{retenirbox}
Pour tout réel \(a\), sur \(]0,+\infty[\) :
\[
(x^a)'=ax^{a-1}.
\]
\end{retenirbox}

\begin{exercicebox}[Puissances]
Dériver :
\[
f(x)=x^{\sqrt2}.
\]
\end{exercicebox}

\begin{correctionbox}
La fonction est définie et dérivable sur \(]0,+\infty[\).  
On applique la formule :
\[
(x^a)'=ax^{a-1}.
\]
Avec :
\[
a=\sqrt2,
\]
on obtient :
\[
f'(x)=\sqrt2\,x^{\sqrt2-1}.
\]
\end{correctionbox}

%=========================================================
\section{Fonctions trigonométriques}

\subsection{Sinus et cosinus}

\begin{theorembox}[Dérivées usuelles]
Les fonctions sinus et cosinus sont dérivables sur \(\R\), et :
\[
(\sin x)'=\cos x,
\]
\[
(\cos x)'=-\sin x.
\]
\end{theorembox}

\begin{theorembox}[Identité fondamentale]
Pour tout \(x\in\R\),
\[
\cos^2 x+\sin^2 x=1.
\]
\end{theorembox}

\subsection{Tangente}

\begin{definitionbox}[Tangente]
La fonction tangente est définie par :
\[
\tan x=\frac{\sin x}{\cos x},
\]
pour :
\[
x\in\R\setminus\left\{\frac{\pi}{2}+k\pi:\ k\in\Z\right\}.
\]
\end{definitionbox}

\begin{theorembox}[Dérivée de la tangente]
La fonction tangente est dérivable sur chaque intervalle de son domaine, et :
\[
(\tan x)'=1+\tan^2 x=\frac1{\cos^2 x}.
\]
\end{theorembox}

\begin{proof}
On utilise la formule du quotient :
\[
(\tan x)'=
\left(\frac{\sin x}{\cos x}\right)'
=
\frac{\cos x\cos x-\sin x(-\sin x)}{\cos^2 x}.
\]
Donc :
\[
(\tan x)'=\frac{\cos^2 x+\sin^2 x}{\cos^2 x}
=
\frac1{\cos^2 x}.
\]
Comme :
\[
1+\tan^2 x
=
1+\frac{\sin^2 x}{\cos^2 x}
=
\frac{\cos^2 x+\sin^2 x}{\cos^2 x}
=
\frac1{\cos^2 x},
\]
on a aussi :
\[
(\tan x)'=1+\tan^2 x.
\]
\end{proof}

\begin{methodbox}[Dériver une fonction trigonométrique composée]
Pour dériver :
\[
\sin(u(x)),\quad \cos(u(x)),\quad \tan(u(x)),
\]
on utilise :
\[
(\sin u)'=u'\cos u,
\]
\[
(\cos u)'=-u'\sin u,
\]
\[
(\tan u)'=\frac{u'}{\cos^2 u}.
\]
\end{methodbox}

\begin{exercicebox}[Trigonométrie]
Dériver :
\[
f(x)=\sin(x^2)\cos x.
\]
\end{exercicebox}

\begin{correctionbox}
On utilise la formule du produit :
\[
f'(x)=(\sin(x^2))'\cos x+\sin(x^2)(\cos x)'.
\]
Or :
\[
(\sin(x^2))'=2x\cos(x^2),
\]
et :
\[
(\cos x)'=-\sin x.
\]
Donc :
\[
f'(x)=2x\cos(x^2)\cos x-\sin(x^2)\sin x.
\]
\end{correctionbox}

%=========================================================
\section{Fonctions trigonométriques réciproques}

\subsection{Arcsinus}

\begin{definitionbox}[Arcsinus]
La fonction sinus réalise une bijection de :
\[
\left[-\frac{\pi}{2},\frac{\pi}{2}\right]
\]
sur :
\[
[-1,1].
\]
Sa réciproque est appelée arcsinus et notée :
\[
\arcsin:[-1,1]\to\left[-\frac{\pi}{2},\frac{\pi}{2}\right].
\]
\end{definitionbox}

\begin{theorembox}[Dérivée de l'arcsinus]
La fonction \(\arcsin\) est dérivable sur \(]-1,1[\), et :
\[
(\arcsin x)'=\frac1{\sqrt{1-x^2}}.
\]
\end{theorembox}

\begin{proof}
Soit :
\[
y=\arcsin x.
\]
Alors :
\[
x=\sin y,
\]
avec :
\[
y\in\left]-\frac{\pi}{2},\frac{\pi}{2}\right[.
\]
Par dérivation d'une réciproque :
\[
(\arcsin)'(x)=\frac1{\cos y}.
\]
Or :
\[
\cos y>0
\]
sur cet intervalle, et :
\[
\cos y=\sqrt{1-\sin^2 y}=\sqrt{1-x^2}.
\]
Donc :
\[
(\arcsin x)'=\frac1{\sqrt{1-x^2}}.
\]
\end{proof}

\subsection{Arccosinus}

\begin{definitionbox}[Arccosinus]
La fonction cosinus réalise une bijection de :
\[
[0,\pi]
\]
sur :
\[
[-1,1].
\]
Sa réciproque est appelée arccosinus :
\[
\arccos:[-1,1]\to[0,\pi].
\]
\end{definitionbox}

\begin{theorembox}[Dérivée de l'arccosinus]
La fonction \(\arccos\) est dérivable sur \(]-1,1[\), et :
\[
(\arccos x)'=-\frac1{\sqrt{1-x^2}}.
\]
\end{theorembox}

\begin{proof}
Soit :
\[
y=\arccos x.
\]
Alors :
\[
x=\cos y,
\]
avec :
\[
y\in]0,\pi[.
\]
Par dérivation d'une réciproque :
\[
(\arccos)'(x)=\frac1{-\sin y}.
\]
Or :
\[
\sin y>0
\]
sur \(]0,\pi[\), et :
\[
\sin y=\sqrt{1-\cos^2 y}=\sqrt{1-x^2}.
\]
Donc :
\[
(\arccos x)'=-\frac1{\sqrt{1-x^2}}.
\]
\end{proof}

\subsection{Arctangente}

\begin{definitionbox}[Arctangente]
La fonction tangente réalise une bijection de :
\[
\left]-\frac{\pi}{2},\frac{\pi}{2}\right[
\]
sur :
\[
\R.
\]
Sa réciproque est appelée arctangente :
\[
\arctan:\R\to\left]-\frac{\pi}{2},\frac{\pi}{2}\right[.
\]
\end{definitionbox}

\begin{theorembox}[Dérivée de l'arctangente]
La fonction \(\arctan\) est dérivable sur \(\R\), et :
\[
(\arctan x)'=\frac1{1+x^2}.
\]
\end{theorembox}

\begin{proof}
Soit :
\[
y=\arctan x.
\]
Alors :
\[
x=\tan y.
\]
Par dérivation d'une réciproque :
\[
(\arctan)'(x)=\frac1{1+\tan^2 y}.
\]
Comme :
\[
\tan y=x,
\]
on obtient :
\[
(\arctan x)'=\frac1{1+x^2}.
\]
\end{proof}

\begin{retenirbox}
Formules à connaître :
\[
(\arcsin x)'=\frac1{\sqrt{1-x^2}},
\]
\[
(\arccos x)'=-\frac1{\sqrt{1-x^2}},
\]
\[
(\arctan x)'=\frac1{1+x^2}.
\]
\end{retenirbox}

%=========================================================
\section{Fonctions hyperboliques}

\subsection{Définitions}

\begin{definitionbox}[Fonctions hyperboliques]
Pour tout \(x\in\R\), on définit :
\[
\ch x=\frac{\e^x+\e^{-x}}2,
\]
\[
\sh x=\frac{\e^x-\e^{-x}}2,
\]
\[
\th x=\frac{\sh x}{\ch x}.
\]
\end{definitionbox}

\begin{theorembox}[Identité fondamentale]
Pour tout \(x\in\R\),
\[
\ch^2 x-\sh^2 x=1.
\]
\end{theorembox}

\begin{proof}
On calcule :
\[
\ch^2 x-\sh^2 x
=
\left(\frac{\e^x+\e^{-x}}2\right)^2
-
\left(\frac{\e^x-\e^{-x}}2\right)^2.
\]
On utilise :
\[
(a+b)^2-(a-b)^2=4ab.
\]
Avec :
\[
a=\e^x,\qquad b=\e^{-x},
\]
on obtient :
\[
\ch^2 x-\sh^2 x
=
\frac{4\e^x\e^{-x}}4
=
1.
\]
\end{proof}

\subsection{Dérivées}

\begin{theorembox}[Dérivées hyperboliques]
Les fonctions \(\ch\), \(\sh\) et \(\th\) sont dérivables sur \(\R\), et :
\[
(\ch x)'=\sh x,
\]
\[
(\sh x)'=\ch x,
\]
\[
(\th x)'=1-\th^2 x=\frac1{\ch^2 x}.
\]
\end{theorembox}

\begin{proof}
Les deux premières formules se déduisent directement des définitions et de :
\[
(\e^x)'=\e^x,
\qquad
(\e^{-x})'=-\e^{-x}.
\]
Pour la tangente hyperbolique :
\[
\th x=\frac{\sh x}{\ch x}.
\]
Donc :
\[
(\th x)'=
\frac{\ch x\cdot\ch x-\sh x\cdot\sh x}{\ch^2 x}
=
\frac{\ch^2 x-\sh^2 x}{\ch^2 x}
=
\frac1{\ch^2 x}.
\]
Comme :
\[
1-\th^2 x
=
1-\frac{\sh^2 x}{\ch^2 x}
=
\frac{\ch^2 x-\sh^2 x}{\ch^2 x}
=
\frac1{\ch^2 x},
\]
on obtient aussi :
\[
(\th x)'=1-\th^2 x.
\]
\end{proof}

%=========================================================
\newpage
\section{Exercices progressifs}

\subsection*{Thème 1 -- Dérivabilité et tangente}

\begin{exercicebox}[1]
Montrer avec la définition que la fonction :
\[
f(x)=x^2
\]
est dérivable en \(a\in\R\), et déterminer \(f'(a)\).
\end{exercicebox}

\begin{exercicebox}[2]
Étudier la dérivabilité en \(0\) de :
\[
f(x)=x|x|.
\]
\end{exercicebox}

\begin{exercicebox}[3]
Déterminer la tangente à la courbe de :
\[
f(x)=\ln x
\]
au point d'abscisse \(1\).
\end{exercicebox}

\subsection*{Thème 2 -- Calculs de dérivées}

\begin{exercicebox}[4]
Dériver :
\[
f(x)=\frac{x^3+1}{x^2+1}.
\]
\end{exercicebox}

\begin{exercicebox}[5]
Dériver :
\[
f(x)=\ln(1+\e^x).
\]
\end{exercicebox}

\begin{exercicebox}[6]
Dériver :
\[
f(x)=\arctan(x^2).
\]
\end{exercicebox}

\subsection*{Thème 3 -- Rolle et accroissements finis}

\begin{exercicebox}[7]
Montrer que l'équation :
\[
x+\sin x=0
\]
admet une unique solution réelle.
\end{exercicebox}

\begin{exercicebox}[8]
Montrer que pour tous \(x,y\in\R\),
\[
|\e^x-\e^y|\leq \e^{\max(x,y)}|x-y|.
\]
\end{exercicebox}

\begin{exercicebox}[9]
Montrer que :
\[
\forall x>0,\qquad \ln x\leq x-1.
\]
\end{exercicebox}

\subsection*{Thème 4 -- Fonctions usuelles}

\begin{exercicebox}[10]
Étudier les variations de :
\[
f(x)=x-\ln x
\]
sur \(]0,+\infty[\).
\end{exercicebox}

\begin{exercicebox}[11]
Montrer que :
\[
\arcsin x+\arccos x=\frac{\pi}{2}
\]
pour tout \(x\in[-1,1]\).
\end{exercicebox}

\begin{exercicebox}[12]
Montrer que :
\[
\ch x\geq1
\]
pour tout \(x\in\R\).
\end{exercicebox}

%=========================================================
\newpage
\section{Corrigés détaillés des exercices progressifs}

\begin{correctionbox}[1]
Soit :
\[
f(x)=x^2.
\]
On calcule le taux d'accroissement en \(a\) :
\[
\frac{f(a+h)-f(a)}{h}
=
\frac{(a+h)^2-a^2}{h}.
\]
On développe :
\[
(a+h)^2-a^2=2ah+h^2.
\]
Donc, pour \(h\neq0\),
\[
\frac{(a+h)^2-a^2}{h}=2a+h.
\]
En faisant tendre \(h\) vers \(0\), on obtient :
\[
f'(a)=2a.
\]
\end{correctionbox}

\begin{correctionbox}[2]
On considère :
\[
f(x)=x|x|.
\]
Pour \(x>0\),
\[
f(x)=x^2.
\]
Pour \(x<0\),
\[
f(x)=-x^2.
\]
On calcule en \(0\) :
\[
\frac{f(h)-f(0)}{h}=\frac{h|h|}{h}=|h|.
\]
Comme :
\[
|h|\to0
\]
quand \(h\to0\), la fonction est dérivable en \(0\), et :
\[
f'(0)=0.
\]
\end{correctionbox}

\begin{correctionbox}[3]
On a :
\[
f(x)=\ln x.
\]
Alors :
\[
f(1)=0,
\]
et :
\[
f'(x)=\frac1x.
\]
Donc :
\[
f'(1)=1.
\]
L'équation de la tangente est :
\[
y=f(1)+f'(1)(x-1).
\]
Donc :
\[
y=x-1.
\]
\end{correctionbox}

\begin{correctionbox}[4]
On pose :
\[
u(x)=x^3+1,
\qquad
v(x)=x^2+1.
\]
Alors :
\[
u'(x)=3x^2,
\qquad
v'(x)=2x.
\]
Comme :
\[
v(x)>0
\]
pour tout \(x\), la fonction est dérivable sur \(\R\), et :
\[
f'(x)=\frac{3x^2(x^2+1)-(x^3+1)2x}{(x^2+1)^2}.
\]
On développe :
\[
f'(x)=\frac{3x^4+3x^2-2x^4-2x}{(x^2+1)^2}.
\]
Donc :
\[
f'(x)=\frac{x^4+3x^2-2x}{(x^2+1)^2}.
\]
\end{correctionbox}

\begin{correctionbox}[5]
On pose :
\[
f(x)=\ln(1+\e^x).
\]
Soit :
\[
u(x)=1+\e^x.
\]
Alors :
\[
u'(x)=\e^x.
\]
Comme :
\[
u(x)>0,
\]
on a :
\[
f'(x)=\frac{u'(x)}{u(x)}
=
\frac{\e^x}{1+\e^x}.
\]
\end{correctionbox}

\begin{correctionbox}[6]
On pose :
\[
f(x)=\arctan(x^2).
\]
La dérivée de \(\arctan u\) est :
\[
\frac{u'}{1+u^2}.
\]
Ici :
\[
u(x)=x^2,
\qquad
u'(x)=2x.
\]
Donc :
\[
f'(x)=\frac{2x}{1+x^4}.
\]
\end{correctionbox}

\begin{correctionbox}[7]
On pose :
\[
f(x)=x+\sin x.
\]
La fonction est dérivable sur \(\R\), et :
\[
f'(x)=1+\cos x.
\]
Comme :
\[
\cos x\geq -1,
\]
on a :
\[
f'(x)\geq0.
\]
Donc \(f\) est croissante.

Pour montrer l'unicité plus précisément, supposons que \(f(a)=f(b)=0\) avec \(a<b\).  
Par Rolle, il existe \(c\in]a,b[\) tel que :
\[
f'(c)=0.
\]
Donc :
\[
1+\cos c=0,
\]
c'est-à-dire :
\[
\cos c=-1.
\]
Cela impose :
\[
c=\pi+2k\pi.
\]
Mais autour d'un tel point, la fonction \(x+\sin x\) reste strictement croissante au sens global.  
Plus simplement, pour \(x<y\),
\[
f(y)-f(x)=y-x+\sin y-\sin x.
\]
Par l'inégalité des accroissements finis :
\[
|\sin y-\sin x|\leq |y-x|=y-x.
\]
Donc :
\[
f(y)-f(x)\geq0.
\]
La fonction est croissante.

On remarque directement que :
\[
f(0)=0.
\]
Si \(x>0\), alors :
\[
x+\sin x>0
\]
sauf cas impossible de compensation totale permanente ; et si \(x<0\), par imparité :
\[
f(x)=-f(-x)<0.
\]
Ainsi l'unique solution est :
\[
x=0.
\]
\end{correctionbox}

\begin{correctionbox}[8]
Soient \(x,y\in\R\).  
Supposons par exemple \(x<y\).  
On applique le théorème des accroissements finis à :
\[
t\mapsto \e^t
\]
sur \([x,y]\).  
Il existe \(c\in]x,y[\) tel que :
\[
\e^y-\e^x=\e^c(y-x).
\]
Comme :
\[
c\leq y=\max(x,y),
\]
on a :
\[
\e^c\leq \e^{\max(x,y)}.
\]
Donc :
\[
|\e^y-\e^x|\leq \e^{\max(x,y)}|y-x|.
\]
Le cas \(y<x\) est analogue.
\end{correctionbox}

\begin{correctionbox}[9]
On pose :
\[
f(x)=x-1-\ln x
\]
sur \(]0,+\infty[\).  
On veut montrer :
\[
f(x)\geq0.
\]
On dérive :
\[
f'(x)=1-\frac1x=\frac{x-1}{x}.
\]
Donc :
\[
f'(x)<0
\quad\text{sur } ]0,1[,
\]
et :
\[
f'(x)>0
\quad\text{sur } ]1,+\infty[.
\]
La fonction \(f\) admet donc un minimum en \(x=1\).  
Or :
\[
f(1)=1-1-\ln1=0.
\]
Ainsi :
\[
f(x)\geq0
\]
pour tout \(x>0\).  
Donc :
\[
\ln x\leq x-1.
\]
\end{correctionbox}

\begin{correctionbox}[10]
On considère :
\[
f(x)=x-\ln x
\]
sur \(]0,+\infty[\).  
On dérive :
\[
f'(x)=1-\frac1x=\frac{x-1}{x}.
\]
Comme \(x>0\), le signe de \(f'(x)\) est celui de \(x-1\).  
Donc :
\[
f'(x)<0\quad\text{sur } ]0,1[,
\]
\[
f'(1)=0,
\]
\[
f'(x)>0\quad\text{sur } ]1,+\infty[.
\]
Ainsi \(f\) est décroissante sur \(]0,1]\), puis croissante sur \([1,+\infty[\).  
Elle admet un minimum en \(1\), de valeur :
\[
f(1)=1.
\]
\end{correctionbox}

\begin{correctionbox}[11]
Soit :
\[
F(x)=\arcsin x+\arccos x.
\]
Sur \(]-1,1[\), on dérive :
\[
F'(x)=\frac1{\sqrt{1-x^2}}-\frac1{\sqrt{1-x^2}}=0.
\]
Donc \(F\) est constante sur \(]-1,1[\).  
En évaluant en \(0\) :
\[
F(0)=\arcsin 0+\arccos 0=0+\frac{\pi}{2}=\frac{\pi}{2}.
\]
Par continuité aux bornes, l'égalité reste vraie sur \([-1,1]\).  
Donc :
\[
\arcsin x+\arccos x=\frac{\pi}{2}.
\]
\end{correctionbox}

\begin{correctionbox}[12]
On sait que :
\[
\ch x=\frac{\e^x+\e^{-x}}2.
\]
Comme :
\[
\e^x>0
\quad\text{et}\quad
\e^{-x}>0,
\]
on peut appliquer l'inégalité arithmético-géométrique :
\[
\frac{\e^x+\e^{-x}}2\geq \sqrt{\e^x\e^{-x}}.
\]
Or :
\[
\sqrt{\e^x\e^{-x}}=\sqrt{1}=1.
\]
Donc :
\[
\ch x\geq1.
\]
\end{correctionbox}

%=========================================================
\newpage
\section{Grand devoir bilan final}

\begin{center}
\Large\textbf{Devoir bilan -- Dérivation et fonctions usuelles}
\end{center}

\subsection*{Exercice 1 -- Questions de cours}

\begin{enumerate}
    \item Définir la dérivabilité d'une fonction en un point.
    \item Montrer que la dérivabilité implique la continuité.
    \item Énoncer le théorème de Rolle.
    \item Énoncer le théorème des accroissements finis.
    \item Énoncer l'inégalité des accroissements finis.
    \item Donner les dérivées de \(\ln x\), \(\e^x\), \(x^a\), \(\arctan x\).
\end{enumerate}

\subsection*{Exercice 2 -- Calcul différentiel}

On considère :
\[
f(x)=\frac{\ln(1+x^2)}{1+x}
\]
sur son domaine de définition.

\begin{enumerate}
    \item Déterminer son domaine de définition.
    \item Calculer \(f'(x)\).
    \item Étudier la dérivabilité éventuelle en \(0\).
\end{enumerate}

\subsection*{Exercice 3 -- Accroissements finis}

Montrer que :
\[
\forall x> -1,\qquad \ln(1+x)\leq x.
\]

\subsection*{Exercice 4 -- Fonction réciproque}

Soit :
\[
f(x)=x+\e^x.
\]

\begin{enumerate}
    \item Montrer que \(f\) réalise une bijection de \(\R\) sur \(\R\).
    \item Montrer que sa réciproque est dérivable.
    \item Donner une formule pour \((f^{-1})'(y)\).
\end{enumerate}

\subsection*{Exercice 5 -- Étude complète}

On considère :
\[
g(x)=x-\arctan x.
\]

\begin{enumerate}
    \item Déterminer le domaine de définition.
    \item Étudier la parité de \(g\).
    \item Calculer \(g'(x)\).
    \item Étudier les variations de \(g\).
    \item Montrer que :
    \[
    \forall x\geq0,\quad \arctan x\leq x.
    \]
\end{enumerate}

\newpage

%=========================================================
\section{Correction détaillée du devoir bilan}

\subsection*{Correction de l'exercice 1}

\textbf{1.}  
La fonction \(f\) est dérivable en \(a\) si la limite :
\[
\lim_{h\to0}\frac{f(a+h)-f(a)}{h}
\]
existe et est finie.  
Cette limite est alors notée :
\[
f'(a).
\]

\textbf{2.}  
Si \(f\) est dérivable en \(a\), alors :
\[
f(a+h)=f(a)+hf'(a)+h\varepsilon(h),
\]
avec :
\[
\varepsilon(h)\to0.
\]
Donc :
\[
f(a+h)-f(a)=hf'(a)+h\varepsilon(h)\to0.
\]
Ainsi :
\[
f(a+h)\to f(a),
\]
donc \(f\) est continue en \(a\).

\textbf{3.}  
Théorème de Rolle : si \(f\) est continue sur \([a,b]\), dérivable sur \(]a,b[\), et si :
\[
f(a)=f(b),
\]
alors il existe \(c\in]a,b[\) tel que :
\[
f'(c)=0.
\]

\textbf{4.}  
Théorème des accroissements finis : si \(f\) est continue sur \([a,b]\) et dérivable sur \(]a,b[\), alors il existe \(c\in]a,b[\) tel que :
\[
f(b)-f(a)=f'(c)(b-a).
\]

\textbf{5.}  
Si \(f\) est dérivable sur un intervalle \(I\) et si :
\[
|f'(x)|\leq K
\]
pour tout \(x\in I\), alors :
\[
|f(x)-f(y)|\leq K|x-y|
\]
pour tous \(x,y\in I\).

\textbf{6.}  
On a :
\[
(\ln x)'=\frac1x,
\]
\[
(\e^x)'=\e^x,
\]
\[
(x^a)'=ax^{a-1}
\quad\text{sur } ]0,+\infty[,
\]
\[
(\arctan x)'=\frac1{1+x^2}.
\]

\subsection*{Correction de l'exercice 2}

On considère :
\[
f(x)=\frac{\ln(1+x^2)}{1+x}.
\]

Le logarithme est défini car :
\[
1+x^2>0
\]
pour tout \(x\in\R\).  
Le dénominateur impose :
\[
1+x\neq0.
\]
Donc :
\[
D_f=\R\setminus\{-1\}.
\]

On pose :
\[
u(x)=\ln(1+x^2),
\qquad
v(x)=1+x.
\]
Alors :
\[
u'(x)=\frac{2x}{1+x^2},
\qquad
v'(x)=1.
\]
Donc :
\[
f'(x)=\frac{u'(x)v(x)-u(x)v'(x)}{v(x)^2}.
\]
Ainsi :
\[
f'(x)=
\frac{\frac{2x(1+x)}{1+x^2}-\ln(1+x^2)}
{(1+x)^2}.
\]

En \(0\), la fonction est bien définie.  
Comme elle est quotient de fonctions dérivables avec dénominateur non nul, elle est dérivable en \(0\).  
On peut calculer :
\[
f'(0)=
\frac{0-\ln1}{1}
=
0.
\]

\subsection*{Correction de l'exercice 3}

On veut montrer :
\[
\ln(1+x)\leq x
\]
pour \(x>-1\).

On pose :
\[
F(x)=x-\ln(1+x).
\]
La fonction \(F\) est définie sur \(]-1,+\infty[\).  
On dérive :
\[
F'(x)=1-\frac1{1+x}
=
\frac{x}{1+x}.
\]
Comme \(1+x>0\), le signe de \(F'(x)\) est celui de \(x\).  
Donc \(F\) est décroissante sur \(]-1,0]\) et croissante sur \([0,+\infty[\).  
Elle admet donc un minimum en \(0\).  
Or :
\[
F(0)=0-\ln1=0.
\]
Donc :
\[
F(x)\geq0.
\]
Ainsi :
\[
\ln(1+x)\leq x.
\]

\subsection*{Correction de l'exercice 4}

On considère :
\[
f(x)=x+\e^x.
\]
La fonction est continue et dérivable sur \(\R\).  
On a :
\[
f'(x)=1+\e^x>0.
\]
Donc \(f\) est strictement croissante sur \(\R\).

De plus :
\[
\lim_{x\to-\infty}(x+\e^x)=-\infty,
\]
car \(\e^x\to0\), et :
\[
\lim_{x\to+\infty}(x+\e^x)=+\infty.
\]
Par le théorème de la bijection continue, \(f\) réalise une bijection de \(\R\) sur \(\R\).

Comme :
\[
f'(x)=1+\e^x\neq0
\]
pour tout \(x\), sa réciproque est dérivable sur \(\R\).  
Pour \(y\in\R\), on a :
\[
(f^{-1})'(y)=\frac1{f'(f^{-1}(y))}.
\]
Donc :
\[
(f^{-1})'(y)=\frac1{1+\e^{f^{-1}(y)}}.
\]

\subsection*{Correction de l'exercice 5}

On considère :
\[
g(x)=x-\arctan x.
\]
La fonction est définie sur \(\R\).

Comme \(x\mapsto x\) est impaire et \(\arctan\) est impaire, la différence de deux fonctions impaires est impaire.  
Donc \(g\) est impaire.

On dérive :
\[
g'(x)=1-\frac1{1+x^2}.
\]
Donc :
\[
g'(x)=\frac{x^2}{1+x^2}.
\]
Ainsi :
\[
g'(x)\geq0
\]
pour tout \(x\), avec égalité seulement en \(0\).  
La fonction \(g\) est donc croissante sur \(\R\).

Comme :
\[
g(0)=0-\arctan0=0,
\]
pour \(x\geq0\), on a :
\[
g(x)\geq g(0)=0.
\]
Donc :
\[
x-\arctan x\geq0.
\]
Ainsi :
\[
\arctan x\leq x.
\]

%=========================================================
\newpage
\section{Fiche de synthèse finale}

\subsection*{1. Dérivabilité}

\[
f'(a)=\lim_{h\to0}\frac{f(a+h)-f(a)}{h}.
\]

Si \(f\) est dérivable en \(a\), alors :
\[
f(a+h)=f(a)+hf'(a)+o(h).
\]

La tangente au graphe de \(f\) en \(a\) est :
\[
y=f(a)+f'(a)(x-a).
\]

\subsection*{2. Dérivabilité et continuité}

\[
f\text{ dérivable en }a
\implique
f\text{ continue en }a.
\]

La réciproque est fausse.

\subsection*{3. Opérations}

\[
(f+g)'=f'+g',
\]
\[
(fg)'=f'g+fg',
\]
\[
\left(\frac fg\right)'=\frac{f'g-fg'}{g^2}.
\]

\[
(f\circ g)'=(f'\circ g)g'.
\]

\subsection*{4. Réciproque}

Si \(f\) est bijective, dérivable, et \(f'\neq0\), alors :
\[
(f^{-1})'(y)=\frac1{f'(f^{-1}(y))}.
\]

\subsection*{5. Extrema}

Si \(f\) admet un extremum local en un point intérieur \(a\), et si \(f\) est dérivable en \(a\), alors :
\[
f'(a)=0.
\]

\subsection*{6. Rolle}

Si \(f\) est continue sur \([a,b]\), dérivable sur \(]a,b[\), et si :
\[
f(a)=f(b),
\]
alors :
\[
\exists c\in]a,b[,\quad f'(c)=0.
\]

\subsection*{7. Accroissements finis}

Si \(f\) est continue sur \([a,b]\) et dérivable sur \(]a,b[\), alors :
\[
\exists c\in]a,b[,\quad
f(b)-f(a)=f'(c)(b-a).
\]

\subsection*{8. Inégalité des accroissements finis}

Si :
\[
|f'(x)|\leq K
\]
sur un intervalle \(I\), alors :
\[
|f(x)-f(y)|\leq K|x-y|.
\]

\subsection*{9. Fonctions usuelles}

\[
(\ln x)'=\frac1x.
\]

\[
(\e^x)'=\e^x.
\]

\[
(x^a)'=ax^{a-1}.
\]

\[
(\sin x)'=\cos x.
\]

\[
(\cos x)'=-\sin x.
\]

\[
(\tan x)'=\frac1{\cos^2 x}=1+\tan^2 x.
\]

\[
(\arcsin x)'=\frac1{\sqrt{1-x^2}}.
\]

\[
(\arccos x)'=-\frac1{\sqrt{1-x^2}}.
\]

\[
(\arctan x)'=\frac1{1+x^2}.
\]

\[
(\ch x)'=\sh x,
\qquad
(\sh x)'=\ch x,
\qquad
(\th x)'=\frac1{\ch^2 x}.
\]

\subsection*{10. Méthodes essentielles}

Pour étudier une fonction :
\begin{enumerate}
    \item déterminer le domaine ;
    \item calculer la dérivée ;
    \item étudier le signe de la dérivée ;
    \item dresser le tableau de variations ;
    \item utiliser les limites aux bornes ;
    \item appliquer le TVI ou le théorème de la bijection si nécessaire.
\end{enumerate}

\begin{center}
\[
\boxed{
\text{La dérivée transforme l'étude locale en information globale sur les variations.}
}
\]
\end{center}

%=========================================================
\section*{Conclusion pédagogique}
\addcontentsline{toc}{section}{Conclusion pédagogique}

Ce chapitre est un pivot de l'analyse réelle.  
La notion de dérivée donne une approximation locale de la fonction.  
Les théorèmes de Rolle et des accroissements finis permettent d'en tirer des conséquences globales.  
Les fonctions usuelles fournissent ensuite un répertoire indispensable pour tous les chapitres suivants : développements limités, intégration, équations différentielles, séries et calcul différentiel.

\begin{center}
\[
\boxed{
\text{Comprendre la dérivation, c'est comprendre comment une fonction varie.}
}
\]
\end{center}

\end{document}